% !PRINTED 499
\weberpage{499}
%<<CONT
cienten von $\varphi'$ sind symmetrische Functionen der Wurzeln der Gleichung
\[ \frac{f(x)}{x - \Theta} = 0, \]
die ihrerseits in $\Omega$ enthalten sind.

Aus der Definition ergiebt sich, dass die Norm eines Productes gleich dem Producte der Normen der Factoren ist, dass also, wenn $\varphi, \psi$ Functionen in $\Omega$ sind,
\weq{7}{N(\varphi\psi) = N(\varphi)\,N(\psi)}
ist.

\begin{center}
\S.\,136.

\gesperrt{Die Functionale eines algebraischen K\"orpers $\Omega$ und der erweiterte K\"orper $\overline{\Omega}$.}
\end{center}

Die Variablen, die in der Theorie der algebraischen Zahlen verwendet werden, haben nicht die Bedeutung von Zeichen f\"ur ver\"anderliche Zahlenreihen, wie man es aus der Functionentheorie gew\"ohnt ist, sondern sie sind lediglich Rechnungssymbole ohne eine selbst\"andige Bedeutung. Bei den Functionen dieser Variablen kommt es eigentlich nur auf die Co\"efficientensysteme an, und die Variablen werden nur dazu benutzt, um die bekannten und gel\"aufigen Regeln der Buchstabenrechnung auf diese Co\"efficientensysteme anzuwenden. Es ist damit freilich nicht ausgeschlossen, dass gelegentlich auch die Zahlen betrachtet werden, die man erh\"alt, wenn man die Variablen durch gewisse Zahlen, z. B. durch rationale Zahlen, ersetzt.

\gesperrt{Demnach f\"uhren wir in unsere Betrachtungen sowohl ganze als gebrochene Functionen von beliebig vielen Ver\"anderlichen ein, deren Co\"efficienten Zahlen eines algebraischen K\"orpers $\Omega$ sind, und setzen fest, dass mit diesen Functionen so gerechnet wird, wie es die Buchstabenrechnung vorschreibt.}

Jede solche Function $\omega$ kann als Quotient zweier ganzer Functionen in $\Omega$,
\weq{1}{\omega = \frac{\varphi}{\psi},}
dargestellt werden, wobei $\psi$ immer von Null verschieden angenommen werden muss. Zwei solche Functionen sind nur dann einander gleich:
\[ \frac{\varphi}{\psi} = \frac{\varphi_1}{\psi_1}, \]
%CONT>>
